\begin{exercise}{20}
Let $X$ be a simplicial set and let $n\le m$. Show that
$\mathrm{sk}_n(\mathrm{sk}_m(X))=\mathrm{sk}_n(X)
=\mathrm{sk}_m(\mathrm{sk}_n(X))$. Deduce that
$\mathrm{cosk}_n(\mathrm{cosk}_m(X))\cong
\mathrm{cosk}_m(\mathrm{cosk}_n(X))$. Is it also true that
$\mathrm{sk}_n(\mathrm{cosk}_m(X))\cong
\mathrm{cosk}_m(\mathrm{sk}_n(X))$? (If not: Provide a counter
example.) Is there a preferred map between these two simplicial sets?

\begin{proof}
Write $S_r=\mathrm{sk}_r$ and $C_r=\mathrm{cosk}_r$.
By Lemma~1.1.39(1), $S_rX$ is the sub-simplicial set generated by
the simplices of $X$ in degrees at most $r$.
Thus $S_nS_mX=S_nX$, since $S_mX$ agrees with $X$ through degree
$n$; and $S_mS_nX=S_nX$, since $S_nX$ is already generated in
degrees at most $n\le m$.

By Lemma~1.1.39(2), the right adjoints of $S_nS_m$, $S_n$, and
$S_mS_n$ are $C_mC_n$, $C_n$, and $C_nC_m$, respectively.
Uniqueness of right adjoints (Lemma~1.1.19) therefore gives canonical
isomorphisms
\[
 C_nC_mX\cong C_nX\cong C_mC_nX.
\]

The mixed isomorphism fails. Take $n=m=0$ and
$X=\Delta^0\amalg\Delta^0$.
The simplicial set $S_0C_0X$ has only its two vertices and their
degeneracies, so it has two $1$-simplices.
In contrast, by Definition~1.1.38,
\[
 (C_0S_0X)_1
 =\operatorname{Hom}(S_0\Delta^1,S_0X)
 \cong X_0\times X_0,
\]
which has four elements. The two simplicial sets cannot be isomorphic.

There is a preferred map in the direction $S_nC_mX\to C_mS_nX$.
Let $u_Z^m\colon Z\to C_mZ$ denote the canonical map.
Definition~1.1.38 shows that $u_X^m$ is a bijection through degree
$m$. Since $S_n$ depends only on degrees at most $n\le m$,
$S_nu_X^m$ is an isomorphism. The preferred comparison is
\[
 S_nC_mX\xrightarrow[\cong]{(S_nu_X^m)^{-1}}S_nX
          \xrightarrow{u_{S_nX}^m}C_mS_nX.
\]
Both the canonical map $u^m$ and the displayed inverse isomorphism
are natural in $X$, so this comparison is natural.
\end{proof}
\end{exercise}
